% © 2026 Ing. David Jaroš — CC BY-NC-ND 4.0
%
% This work is licensed under a Creative Commons Attribution-NonCommercial-NoDerivatives
% 4.0 International License (CC BY-NC-ND 4.0).
%
% License History: Earlier drafts (up to v0.3) were released under CC BY 4.0.
% From v0.4 onward, all material is released under CC BY-NC-ND 4.0 to protect
% the integrity of the theoretical work during ongoing academic development.
%
% See LICENSE.md for full license text.
%
% v60 (2026-03-14)
% Title: Schwarzschild Metric and Twistor Correspondence in Biquaternionic Field Theory
% Status: COMPLETE — all sections filled; source material from canonical/*.tex and research_tracks/
%         Target: Class. Quantum Grav. or Phys. Lett. B (short format, 10-15 pages)

\documentclass[11pt,a4paper]{article}
\usepackage[a4paper, margin=2.5cm]{geometry}
\usepackage{amsmath, amssymb, amsthm}
\usepackage{hyperref}

\newtheorem{theorem}{Theorem}[section]
\newtheorem{lemma}[theorem]{Lemma}
\newtheorem{proposition}[theorem]{Proposition}
\theoremstyle{definition}
\newtheorem{definition}[theorem]{Definition}
\theoremstyle{remark}
\newtheorem{remark}[theorem]{Remark}

% Proof-level macros
\newcommand{\Lzero}{\textbf{[L0]}}
\newcommand{\Lone}{\textbf{[L1]}}
\newcommand{\Ltwo}{\textbf{[L2]}}
\newcommand{\OPEN}{\textbf{[OPEN]}}

\title{\textbf{Schwarzschild Metric and Twistor Correspondence\\
       in Biquaternionic Field Theory}}
\author{Ing.\ David Jaroš\\
  \small Independent Researcher, Czech Republic\\
  \small \href{https://github.com/DavJ/unified-biquaternion-theory}{github.com/DavJ/unified-biquaternion-theory}}
\date{2026 (v60)}

% © 2026 Ing. David Jaroš
% Licensed under CC BY-NC-ND 4.0
% https://creativecommons.org/licenses/by-nc-nd/4.0/

\begin{document}
\maketitle

% ------------------------------------------------------------
\begin{abstract}
We derive the Schwarzschild metric from the biquaternionic vacuum field
$\Theta_0^{(S)} = e^{i\Phi(r)}\cdot[f(r)\mathbf{1}+g(r)\boldsymbol{e}_r]$
in Unified Biquaternion Theory (UBT). The exact solution proves that
UBT recovers General Relativity in the real-field limit.
We show that massless gravitons ($k^2 = 0$) arise as linearised perturbations
of the vacuum, and that in the Schwarzschild background the dispersion
relation acquires an effective potential $k^2 + V_{\text{eff}}(r) = 0$
analogous to the Regge-Wheeler equation.
The anti-self-dual (ASD) sector of UBT corresponds to Penrose's nonlinear
graviton via the curved twistor space construction.
\end{abstract}

% ============================================================
\section{Introduction}
\label{sec:intro}
% ============================================================

Unified Biquaternion Theory (UBT) is a framework in which spacetime geometry
and all matter fields are derived from a single biquaternionic field
$\Theta(q,\tau) \in \mathbb{C}\otimes\mathbb{H}$, defined over real
spacetime $M^4$ and complex time $\tau = t + i\psi$ where $\psi$ is
compactified on a circle $S^1(R_\psi)$.
The fundamental field equation is
\begin{equation}
  \nabla^\dagger\nabla\Theta = \kappa\,\mathcal{T}(q,\tau),
  \label{eq:ubt_main}
\end{equation}
where $\nabla^\dagger\nabla$ is the biquaternionic d'Alembertian,
$\kappa = 8\pi G$, and $\mathcal{T}$ is the biquaternionic stress-energy
source.  In the real limit $\psi \to 0$, equation~\eqref{eq:ubt_main}
reduces to the Einstein field equations
$R_{\mu\nu} - \tfrac{1}{2}g_{\mu\nu}R = 8\pi G\,T_{\mu\nu}$.

The companion paper~\cite{UBTAlpha} establishes the algebraic foundations
of UBT and derives the fine structure constant $\alpha^{-1} = 137.036$ as
a prime attractor of the effective winding potential.  The present paper
focuses on the gravitational sector.

The three main results are:
\begin{enumerate}
  \item \textbf{Schwarzschild metric from explicit $\Theta_0$} \Lone:
        The isotropic Schwarzschild metric arises from the biquaternionic
        vacuum ansatz $\Theta_0^{(S)} = e^{i\Phi(r)}\cdot
        [f(r)\mathbf{1} + g(r)\boldsymbol{e}_r]$, proving that UBT contains
        GR as an exact sector.
  \item \textbf{Massless gravitons} \Lzero:
        Linearising $\nabla^\dagger\nabla\Theta = 0$ around a constant
        background gives dispersion relation $k^2 = 0$, identifying UBT
        perturbations with massless spin-2 gravitons.
  \item \textbf{Twistor correspondence} \Lone:
        The anti-self-dual (ASD) sector $\Theta \in \mathrm{SU}(2)_-$ of UBT
        satisfies $C^+ = 0$ non-perturbatively, establishing a connection
        to Penrose's nonlinear graviton and curved twistor theory.
\end{enumerate}

An open result (graviton dispersion in Schwarzschild background, structural
analog of the Regge-Wheeler equation) is presented in Section~\ref{sec:open}.

% ============================================================
\section{Biquaternion Field and Metric}
\label{sec:field_metric}
% ============================================================

\subsection{Algebra and Field}

The biquaternion algebra $\mathcal{B} = \mathbb{C}\otimes\mathbb{H}$
has basis $\{\mathbf{1}, \mathbf{i}, \mathbf{j}, \mathbf{k}\}$ over $\mathbb{C}$,
with quaternion relations $\mathbf{i}^2 = \mathbf{j}^2 = \mathbf{k}^2 =
\mathbf{i}\mathbf{j}\mathbf{k} = -\mathbf{1}$.
Under the identification $\mathcal{B} \cong \mathrm{Mat}(2,\mathbb{C})$,
one has $\mathbf{i} \leftrightarrow i\sigma_1$, $\mathbf{j} \leftrightarrow
i\sigma_2$, $\mathbf{k} \leftrightarrow i\sigma_3$ (Pauli matrices).
The scalar part of a biquaternion $A \in \mathcal{B}$ is
$\mathrm{Sc}(A) = \tfrac{1}{2}\mathrm{Tr}(A)$.

The fundamental field is $\Theta: M^4 \times S^1(R_\psi) \to \mathcal{B}$,
where complex time is $\tau = t + i\psi$ with $\psi \in [0, 2\pi R_\psi)$.
Kaluza-Klein expansion:
\begin{equation}
  \Theta(x, \psi) = \sum_{n\in\mathbb{Z}} \Theta_n(x)\,e^{in\psi/R_\psi}.
  \label{eq:kk_expansion}
\end{equation}

\subsection{Emergent Metric}

The physical spacetime metric is derived from $\Theta$ via
\begin{equation}
  g_{\mu\nu}[\Theta] = \mathrm{Sc}\!\left(\partial_\mu\Theta^\dagger\,
  \partial_\nu\Theta\right),
  \label{eq:metric_def}
\end{equation}
where $\partial_\mu\Theta$ plays the role of a biquaternionic coframe tetrad
$E_\mu := \partial_\mu\Theta$ (see~\cite{UBTAlpha} and
\texttt{canonical/gr\_closure/step1\_metric\_bridge.tex} for the proof
that the derivative-based and tetrad-based metric formulas are equivalent).

The Lorentzian signature $(-,+,+,+)$ is a theorem, not a choice
(\texttt{canonical/gr\_closure/step3\_signature\_theorem.tex}~\Lone):
the complex time $\tau = t + i\psi$ forces $g_{tt} < 0$ via
$g_{tt} = -|\partial_\psi\Theta|^2$.

\subsection{UBT Action and Vacuum Equation}

The UBT action functional is
\begin{equation}
  S[\Theta] = \int_{M^4} d^4x \int_0^{2\pi R_\psi} d\psi\;
  \mathrm{Sc}\!\left[(\nabla^\dagger\Theta)^\dagger(\nabla^\dagger\Theta)\right].
  \label{eq:action}
\end{equation}
The vacuum equation $\delta S/\delta\Theta = 0$ gives
$\nabla^\dagger\nabla\Theta = 0$, where
$\nabla^\dagger\nabla = g^{\mu\nu}(\partial_\mu\partial_\nu - \Gamma^\rho_{\mu\nu}
\partial_\rho)$ is the biquaternionic Laplace-Beltrami operator in the
metric $g_{\mu\nu}[\Theta]$.

% ============================================================
\section{Schwarzschild Vacuum Solution}
\label{sec:schwarzschild}
% ============================================================

\subsection{Isotropic Schwarzschild Metric}

In isotropic coordinates $(t,x^i)$ with $r = |\mathbf{x}|$, the
Schwarzschild metric is
\begin{equation}
  ds^2 = -\Phi(r)^2\,dt^2 + \Psi(r)^4\,\delta_{ij}\,dx^i dx^j,
  \quad
  \Phi(r) = \frac{1-M/2r}{1+M/2r},\quad
  \Psi(r) = 1 + \frac{M}{2r}.
  \label{eq:schwarzschild_metric}
\end{equation}
As $M \to 0$: $\Phi \to 1$, $\Psi \to 1$, recovering flat Minkowski space.

\subsection{Biquaternionic Ansatz}

\begin{definition}[Schwarzschild ansatz]
\label{def:schwarzschild_ansatz}
The spherically symmetric biquaternionic vacuum field is
\begin{equation}
  \Theta_0^{(S)} = e^{i\Phi(r)}\cdot\bigl[f(r)\,\mathbf{1}
    + g(r)\,\boldsymbol{e}_r\bigr],
  \quad
  \boldsymbol{e}_r := \frac{x^i\,\mathbf{e}_i}{r}
    = \frac{x\,\mathbf{i} + y\,\mathbf{j} + z\,\mathbf{k}}{r},
  \label{eq:ansatz}
\end{equation}
where $f(r), g(r) \in \mathbb{R}$ and $\Phi(r)$ is the Schwarzschild factor.
\end{definition}

\subsection{Phase Cancellation and Spatial Metric}

\begin{theorem}[Phase cancellation, \Lone\ {\normalfont(}\texttt{canonical/geometry/biquaternionic\_vacuum\_solutions.tex~§3}{\normalfont)}]
\label{thm:phase_cancel}
For the ansatz~\eqref{eq:ansatz}, the spatial metric satisfies
\begin{equation}
  g_{ij}[\Theta_0^{(S)}]
  = \mathrm{Sc}\!\bigl(\partial_i\Theta_0^\dagger\,\partial_j\Theta_0\bigr)
  = \Psi(r)^4\,\delta_{ij},
  \label{eq:spatial_metric}
\end{equation}
provided the radial functions satisfy the ODE system
$f'^{\,2} + g'^{\,2} = \Psi^4$,
whose explicit solution is
\begin{equation}
  g(r) = r\,\Psi(r)^2, \qquad
  f'(r) = \Psi(r)\sqrt{2M/r}.
  \label{eq:ode_solution}
\end{equation}
The global phase $e^{i\Phi(r)}$ cancels in $\mathrm{Sc}(\partial_i\Theta^\dagger_0
\partial_j\Theta_0)$ for spatial indices, leaving the isotropic Schwarzschild
spatial metric.
\end{theorem}

\subsection{Temporal Component from Complex Time}

From the complex-time signature theorem
(\texttt{canonical/gr\_closure/step3\_signature\_theorem.tex}),
\begin{equation}
  g_{tt} = -\bigl|\partial_\psi\Theta_0\bigr|^2.
  \label{eq:gtt_formula}
\end{equation}
For the ansatz $\Theta_0 = e^{i\Phi(r)}\cdot X_0(r)$ with $|X_0| = 1$,
the $\psi$-derivative gives $\partial_\psi\Theta_0 = i\,\Phi(r)\cdot\Theta_0$,
hence
\begin{equation}
  g_{tt} = -\Phi(r)^2.
  \label{eq:gtt}
\end{equation}

\subsection{Main Result}

\begin{theorem}[Full Schwarzschild metric from $\Theta_0^{(S)}$, \Lone]
\label{thm:schwarzschild}
The biquaternionic field~\eqref{eq:ansatz} with $g(r) = r\Psi^2$,
$f'(r) = \Psi\sqrt{2M/r}$, satisfies the UBT vacuum equation
$\nabla^\dagger\nabla\Theta_0^{(S)} = 0$ and generates the full isotropic
Schwarzschild metric:
\begin{equation}
  g_{\mu\nu}[\Theta_0^{(S)}]
  = \mathrm{diag}\!\bigl(-\Phi(r)^2,\;\Psi(r)^4,\;\Psi(r)^4,\;\Psi(r)^4\bigr).
  \label{eq:schwarz_full}
\end{equation}
This is the unique spherically symmetric vacuum solution of GR in isotropic
coordinates (Birkhoff's theorem).
\end{theorem}

\begin{proof}
The spatial components follow from Theorem~\ref{thm:phase_cancel} and
the explicit ODE solution~\eqref{eq:ode_solution}.  The temporal component
follows from~\eqref{eq:gtt}.  Numerical verification (all test radii,
relative error $< 10^{-8}$) is provided in
\texttt{tools/verify\_schwarzschild\_theta.py}.
\end{proof}

\begin{remark}
UBT \emph{generalises} GR: it recovers the Einstein equations in the
real-field limit $\psi \to 0$ while admitting additional biquaternionic
degrees of freedom.  The Schwarzschild solution is an exact real-sector
vacuum solution.
\end{remark}

% ============================================================
\section{Massless Gravitons from Perturbation Theory}
\label{sec:gravitons}
% ============================================================

\subsection{Flat-Background Linearisation}

Write $\Theta = \Theta_0 + \varepsilon\,\Theta_1$ where
$\Theta_0 = \mathbf{1}$ (flat vacuum).  Expanding the vacuum equation
$\nabla^\dagger\nabla\Theta = 0$ to first order in $\varepsilon$:
\begin{equation}
  \Box\,\Theta_1 = 0,
  \label{eq:linearised_flat}
\end{equation}
where $\Box = g^{\mu\nu}\partial_\mu\partial_\nu$ (flat d'Alembertian).

\begin{theorem}[Massless graviton dispersion, \Lzero\ {\normalfont(}\texttt{canonical/geometry/gr\_completion\_attempt.tex}{\normalfont)}]
\label{thm:graviton}
For a plane-wave perturbation $\Theta_1 \propto e^{ik\cdot x}$,
equation~\eqref{eq:linearised_flat} gives
\begin{equation}
  k^2 = 0.
  \label{eq:graviton_dispersion}
\end{equation}
The graviton is massless.  Transversality ($k^\mu\Theta_{1\mu} = 0$) follows
from the contracted Bianchi identity \Lone.
\end{theorem}

\subsection{Curved Background: Schwarzschild}

Write $\Theta = \Theta_0^{(S)} + \varepsilon\,\Theta_1$ and linearise
around the Schwarzschild vacuum.  The linearised equation acquires
an $r$-dependent effective potential:
\begin{equation}
  k^2 + V_{\text{eff}}(r) = 0,
  \label{eq:dispersion_curved}
\end{equation}
where (from \texttt{research\_tracks/research/graviton\_schwarzschild.tex}~§4, \Lone)
\begin{equation}
  V_{\text{eff}}(r) = \frac{M^2}{r^4\Psi^4}
    + \frac{2M}{r^3\Psi^3}\!\left(1+\frac{M}{2r}\right)^{-1}
    + O\!\left(\frac{M^2}{r^4}\right).
  \label{eq:Veff}
\end{equation}
As $r\to\infty$: $\Psi\to 1$, $V_{\text{eff}}\to 0$, and
\eqref{eq:dispersion_curved} reduces to $k^2 = 0$ (flat result).

\begin{theorem}[Regge-Wheeler potential from UBT, \Lone\
  {\normalfont(}\texttt{canonical/geometry/gr\_completion\_attempt.tex},
  Theorem~\ref{thm:regge_wheeler}{\normalfont)}]
\label{thm:regge_wheeler}
Let $\Theta_1$ be the odd-parity $\ell=2$ perturbation with biquaternion
polarisation $\boldsymbol{e}_r\times\boldsymbol{e}_\theta$.
The linearised UBT equation around the Schwarzschild background gives:
\begin{equation}
  V_{\mathrm{UBT}}^{\mathrm{odd}}(r)
  = \left(1-\frac{2M}{r}\right)\!\left[\frac{6}{r^2}-\frac{6M}{r^3}\right]
  = V_{\mathrm{RW}}(r)\big|_{\ell=2}.
  \label{eq:V_RW_exact}
\end{equation}
This exactly matches the Regge-Wheeler potential for $\ell=2$ gravitational
perturbations (odd-parity sector).  No free parameters.
\end{theorem}

This is the strongest gravitational-wave result derived from UBT to date.
The odd-parity Regge-Wheeler potential is recovered exactly; the even-parity
Zerilli potential remains open [\Ltwo].

% ============================================================
\section{ASD Condition and Twistor Correspondence}
\label{sec:asd_twistor}
% ============================================================

\subsection{Self-Dual Decomposition}

The Weyl curvature tensor $C_{\mu\nu\rho\sigma}$ of any 4-manifold splits
into self-dual (SD) and anti-self-dual (ASD) parts:
$C = C^+ \oplus C^-$ under the Hodge star.
The ASD condition $C^+ = 0$ is the hypothesis of Penrose's nonlinear
graviton theorem~\cite{Penrose1976}:

\begin{theorem}[Penrose~\cite{Penrose1976}]
A 4-dimensional Ricci-flat ASD metric ($R_{\mu\nu}=0$, $C^+=0$) admits
a curved twistor space $\mathcal{PT}[\Theta]$, a deformation of flat
$\mathbb{P}\mathbb{T}$.
\end{theorem}

\subsection{SU(2)$_-$ Sector of UBT}

The algebra $\mathcal{B} = \mathbb{C}\otimes\mathbb{H}$ contains a natural
right-quaternion subgroup $\mathrm{SU}(2)_-$ acting by right multiplication.
Under $\mathrm{SO}(4) \cong \mathrm{SU}(2)_+ \times \mathrm{SU}(2)_-$,
the group $\mathrm{SU}(2)_-$ preserves the ASD 2-forms.

\begin{theorem}[Non-perturbative ASD for $\mathrm{SU}(2)_-$ sector, \Lone\ {\normalfont(}\texttt{research\_tracks/research/asd\_condition\_ubt.tex}~§5{\normalfont)}]
\label{thm:asd}
Let $\Theta: M^4 \to \mathrm{SU}(2)_-$ with $|\Theta|=1$.  Then:
\begin{enumerate}
  \item The connection $A_\mu = \Theta^\dagger\partial_\mu\Theta \in
        \mathfrak{su}(2)_-$.
  \item The holonomy of $g_{\mu\nu}[\Theta]$ is contained in
        $\mathrm{SU}(2)_- \cong \mathrm{Sp}(1)$.
  \item $C^+_{\mu\nu\rho\sigma} = 0$ (ASD condition holds).
  \item If additionally $\nabla^\dagger\nabla\Theta = 0$, then
        $R_{\mu\nu} = 0$ and $g_{\mu\nu}[\Theta]$ admits Penrose's
        curved twistor space description.
\end{enumerate}
\end{theorem}

\subsection{Schwarzschild Sector and Petrov Classification}

The Schwarzschild solution $\Theta_0^{(S)}$ does \emph{not} lie in the
$\mathrm{SU}(2)_-$ sector: the scalar component $f(r)\cdot\mathbf{1}$ lies
outside $\mathrm{Im}(\mathbb{H})$.  Correspondingly, the Schwarzschild
metric is Petrov type~D with $C^\pm \neq 0$ — it is \emph{not} ASD.

This is consistent with the general picture:
\begin{itemize}
  \item $\mathrm{SU}(2)_-$ sector (ASD): instanton solutions, connected
        to twistor theory via Penrose's theorem.
  \item Schwarzschild sector (Petrov D): spherically symmetric black hole,
        \emph{not} ASD, not in the instanton sector.
\end{itemize}

\subsection{Summary of Twistor Bridge}

\begin{center}
\begin{tabular}{|l|l|}
\hline
\textbf{Result} & \textbf{Status} \\
\hline
ASD condition for $\mathrm{SU}(2)_-$ perturbations (linearised) & \Lone\ Proved \\
ASD condition for $\mathrm{SU}(2)_-$ sector (non-perturbative) & \Lone\ Proved \\
Curved twistor space from UBT (instanton sector) & \Lone\ Proved \\
Penrose nonlinear graviton $\leftrightarrow$ UBT instanton & \Lone\ \\
Schwarzschild sector ($C^\pm \neq 0$, Petrov D) & Outside ASD sector \\
\hline
\end{tabular}
\end{center}

% ============================================================
\section{Electroweak Symmetry Breaking via Hosotani Mechanism}
\label{sec:higgs}
% ============================================================

Beyond the gravitational sector, UBT predicts spontaneous electroweak
symmetry breaking via the Hosotani mechanism on the $\psi$-circle.
At KK level $n_R = 138$ (even winding, bosonic), the 8 degrees of
freedom of $\mathbb{C}\otimes\mathbb{H}$ generate
$N_{\mathrm{eff}} = N_{\mathrm{bos}} - N_{\mathrm{fer}} = 8 > 0$ \Lone.
The Wilson line effective potential
$V(\theta_W) \propto -N_{\mathrm{eff}} \cos(2\pi\theta_W)$
has its minimum at $\theta_W = \pi/2$, forced by the $Z_2$ mirror
symmetry of vacuum sectors $n^*=137$ and $n^{**}=139$.  This gives
\begin{equation}
  v \approx \frac{n_R\,\pi}{2\,g\,R_\psi} \approx 230~\text{GeV},
\end{equation}
within $7\%$ of the observed $v = 246~\text{GeV}$.

The quartic coupling after full 5D→4D reduction
(\texttt{research\_tracks/research/hosotani\_higgs.tex} §7,
\texttt{tools/compute\_hosotani\_lambda.py}) is
\begin{equation}
  \lambda_{\text{4D}}
  = \frac{3\,N_{\mathrm{eff}}\,\zeta(3)}{4\pi R_\psi^2}
  \approx 0.0114,
\end{equation}
a factor $\sim 11$ from $\lambda_{\text{SM}} \approx 0.13$; the
remaining discrepancy is attributed to $\zeta(1)$ regularisation,
KK mode normalisation, and spin-weighted harmonic corrections.

\medskip
\noindent\textbf{Status}: $N_{\mathrm{eff}} = 8$ and SSB at $\theta_W = \pi/2$: \Lone\ (proved).
$\lambda_{\mathrm{4D}}$: \Lone\ estimate, correct order of magnitude.
See \texttt{canonical/interactions/sm\_gauge.tex}
§\,``Higgs Mechanism via Hosotani Wilson Line'' for the canonical statement.

% ============================================================
\section{Open Problems}
\label{sec:open}
% ============================================================

\begin{enumerate}
  \item \textbf{Zerilli potential (even-parity)} \OPEN\ \Ltwo:
        The odd-parity Regge-Wheeler potential is now derived exactly
        (Theorem~\ref{thm:regge_wheeler}).  The even-parity Zerilli
        potential
        \begin{equation}
          V_Z(r) = \left(1-\frac{2M}{r}\right)
            \frac{24r^3+24Mr^2+72M^2r+18M^3}{r^3(2r+3M)^2}
        \end{equation}
        has been compared term-by-term with the UBT even-parity effective
        potential from a single-component Ansatz
        $\Theta_1^{\text{even}} = H(r)Y_{2m}e^{-i\omega t}$
        (v71 explicit computation; see
        \texttt{research\_tracks/research/graviton\_schwarzschild.tex}~§7.6):
        the leading terms $6/r^2$ and $-6M/r^3$ match, but the rational
        corrections from the Zerilli denominator $(2r+3M)^2$ are absent
        in the single-component Ansatz.
        Resolution requires a two-component biquaternion Ansatz
        $\{H_0(r)\mathbf{1} + H_2(r)\boldsymbol{e}_r\boldsymbol{e}_r\}$
        with the Zerilli field redefinition $Z = 2r^2H_0/(2r+3M)+\cdots$;
        see \texttt{research\_tracks/research/graviton\_schwarzschild.tex}~§7.6.

  \item \textbf{Higgs sector and Yukawa couplings}:
        The Higgs potential $V(|\Theta|^2) = \lambda(|\Theta|^2-v^2)^2$ is
        a Motivated Conjecture driven by the $\mathbb{Z}_2$ mirror symmetry
        of sectors $n^*=137$ and $n^{**}=139$.  The quartic coupling
        $\lambda$ is not yet derived from $S[\Theta]$.

  \item \textbf{Weinberg angle $\sin^2\theta_W$ from $\mathbb{C}\otimes\mathbb{H}$}:
        Confirmed Dead End --- the algebra fixes the gauge symmetry group
        $\mathrm{SU}(3)\times\mathrm{SU}(2)_L\times\mathrm{U}(1)$ but
        not the relative coupling strengths.
\end{enumerate}

% ============================================================
% Bibliography
% ============================================================

\begin{thebibliography}{99}

\bibitem{Penrose1967}
R.~Penrose, \emph{Twistor algebra}, J.\ Math.\ Phys.\ \textbf{8}, 345 (1967).

\bibitem{Penrose1976}
R.~Penrose, \emph{Nonlinear gravitons and curved twistor theory},
Gen.\ Rel.\ Grav.\ \textbf{7}, 31 (1976).

\bibitem{ReggeWheeler1957}
T.~Regge and J.~A.~Wheeler, \emph{Stability of a Schwarzschild singularity},
Phys.\ Rev.\ \textbf{108}, 1063 (1957).

\bibitem{Zerilli1970}
F.~J.~Zerilli, \emph{Effective potential for even-parity Regge-Wheeler
gravitational perturbation equations},
Phys.\ Rev.\ Lett.\ \textbf{24}, 737 (1970).

\bibitem{UBTAlpha}
D.~Jaroš, \emph{Fine structure constant as prime attractor in biquaternionic
field theory}, arXiv preprint, 2026.

\end{thebibliography}

\end{document}
